% !TeX root = ../applied-mathematics1-companion.tex
\chapter{Discrete Dynamical Systems}\label{chap:applied-discrete}

\section{Building a discrete population model}
A discrete-time model describes a state at successive time steps:
\[
x_0,\ x_1,\ x_2,\ldots,\qquad x_{t+1}=f(x_t).
\]
It is appropriate, for example, for a population with distinct annual
breeding seasons. The initial condition $x_0$ must be specified.
Population values are treated as real expected sizes or large-population
approximations; the recurrence need not give integer outputs.
A physically meaningful population must satisfy $N_t\geq0$.

\subsection{Constant reproduction rate}
Suppose every individual produces on average $r\geq0$ offspring and the
old generation is replaced completely. Then
\[
N_{t+1}=rN_t,\qquad N_t=r^tN_0.
\]
Repeated substitution or induction proves the solution. If $N_0>0$:
\begin{itemize}
\item $0\leq r<1$: the population tends to zero;
\item $r=1$: the population stays constant;
\item $r>1$: it grows without bound.
\end{itemize}
For $r>0$, $r^t=e^{t\log r}$, hence the description \emph{exponential}
growth or decay. Time is still discrete even when the formula uses an
exponential function.

\subsection{Fixed harvesting}
If $h\geq0$ individuals are removed after reproduction, the algebraic model is
\[
N_{t+1}=rN_t-h.
\]
For $r\ne1$, iteration gives
\begin{align*}
N_t&=r^tN_0-h(1+r+\cdots+r^{t-1})\\
&=r^tN_0-h\frac{r^t-1}{r-1}\\
&=N^*+r^t(N_0-N^*),\qquad N^*=\frac h{r-1}.
\end{align*}
For $r=1$, $N_t=N_0-th$.

If $r>1$ and $h>0$, $N^*>0$ is a threshold:
\begin{itemize}
\item $N_0>N^*$ produces unbounded growth in this model;
\item $N_0=N^*$ produces a constant population;
\item $N_0<N^*$ eventually makes the algebraic formula nonpositive.
\end{itemize}
Negative population predictions indicate failure of the unmodified model.
The 5 October handwritten notes explicitly repair it as
\[
N_{t+1}=\max\{rN_t-h,0\}.
\]
Under this repaired model, a below-threshold population reaches zero in
finite time and remains extinct. The equality case $N_0=N^*$ is an
equilibrium, not extinction. For $h=0$, this reduces to reproduction alone.

\begin{example}[Checking a harvesting threshold]
As an added illustration, let $r=2$ and $h=10$. Then $N^*=10$.
Starting at $12$ gives $N_t=10+2^{t+1}$; starting at $10$ keeps it at $10$.
Starting at $8$ gives $8,6,2,0,0,\ldots$ in the clipped model.
The unclipped next value after $2$ would be $-6$, which is not a population.
\end{example}

\section{The logistic map and nondimensionalisation}
To represent reduced reproduction at high density, take
\[
N_{t+1}=rN_t\left(1-\frac{N_t}{K}\right),\qquad K>0.
\]
Set $x_t=N_t/K$. Dividing by $K$ gives the \emph{logistic map}
\begin{equation}\label{eq:applied-logistic-map}
x_{t+1}=rx_t(1-x_t)=f(x_t).
\end{equation}
Both $x_t$ and $r$ are dimensionless. Scaling removes $K$ from the
recurrence, leaving one parameter controlling the dynamics.

\begin{remark}[What $K$ means here]
In this discrete model, $K$ is the density scale where the next-generation
reproduction factor vanishes. The nonzero equilibrium is
$N^*=K(1-1/r)$, not $K$. This differs from the continuous logistic model,
whose carrying capacity is itself the positive equilibrium.
\end{remark}
For $0\leq x\leq1$, $0\leq x(1-x)\leq1/4$. Consequently, if
$0\leq r\leq4$, the interval $[0,1]$ is invariant: an initial value there
stays there. For $0<r<4$ and $0<x_0<1$, subsequent values remain in $(0,1)$.
For larger $r$, some initial values leave the biologically meaningful range.

\section{Fixed points and cobweb diagrams}
\begin{definition}[Fixed point]
A fixed point of $x_{t+1}=f(x_t)$ is a value $x^*$ with $f(x^*)=x^*$.
Starting exactly at $x^*$ gives a constant sequence. It is also called a
steady state or an equilibrium of the map.
\end{definition}
Fixed points are intersections of $y=f(x)$ with $y=x$.
For the logistic map, assuming $r>0$,
\[
x^*=rx^*(1-x^*)\quad\Longrightarrow\quad
x^*=0\quad\text{or}\quad x^*=1-\frac1r.
\]
The second point is positive only for $r>1$ and coincides with zero at $r=1$.

\studyguide{To construct a cobweb, plot $y=f(x)$ and $y=x$. From $x_0$
on the horizontal axis, go vertically to the curve: the height is $x_1$.
Go horizontally to the diagonal: the horizontal coordinate is now $x_1$.
Repeat vertically to the curve and horizontally to the diagonal.}
% !TeX root = ../applied-mathematics1-companion.tex
% Illustrative attracting affine map; the chapter explains general maps.
\begin{center}
\begin{tikzpicture}[scale=3.7,>=Stealth]
\draw[->] (0,0)--(1.2,0) node[right] {$x$};
\draw[->] (0,0)--(0,1.2) node[above] {$y$};
\draw[muted] (0,0)--(1.1,1.1) node[right] {$y=x$};
\draw[thick,ink] (0,.4)--(1.1,.95) node[right] {$y=f(x)=\tfrac12x+\tfrac25$};
\draw[thick,accent,->] (.1,0)--(.1,.45);
\draw[thick,accent,->] (.1,.45)--(.45,.45);
\draw[thick,accent,->] (.45,.45)--(.45,.625);
\draw[thick,accent,->] (.45,.625)--(.625,.625);
\draw[thick,accent] (.625,.625)--(.625,.7125)--(.7125,.7125)
 --(.7125,.75625)--(.75625,.75625)--(.75625,.778125)--(.778125,.778125);
\fill[ink] (.8,.8) circle (.012) node[above left] {$x^*=4/5$};
\node[below] at (.1,0) {$x_0$};
\end{tikzpicture}

{\small Cobweb steps: vertical to the map, horizontal to the diagonal.}
\end{center}


\section{Taylor approximation and local stability}
For a sufficiently differentiable function near $a$, a finite Taylor expansion
has the form
\[
f(a+h)=f(a)+f'(a)h+\frac12f''(a)h^2+O(h^3).
\]
Here $O(h^3)$ denotes an error bounded in magnitude by a constant times
$|h|^3$ for sufficiently small $h$, under appropriate smoothness assumptions.
In particular $f(a+h)=f(a)+f'(a)h+o(h)$ for a differentiable $f$.
The formal infinite Taylor series does not automatically equal every smooth
function; local approximation is all that is required for linearisation.

\begin{example}[Taylor expansion of $\log x$ near $2$]
Since $f'(x)=1/x$, $f''(x)=-1/x^2$ and $f'''(x)=2/x^3$,
\[
\log x=\log2+\frac{x-2}{2}-\frac{(x-2)^2}{8}
+\frac{(x-2)^3}{24}+O((x-2)^4).
\]
For a first-order approximation, keep only $\log2+(x-2)/2$.
\end{example}

\begin{definition}[Stability and attraction]
An equilibrium is \emph{stable} if arbitrarily small initial perturbations
can be kept arbitrarily small for all future times. It is \emph{locally
asymptotically stable} if it is stable and nearby states also converge to it.
An unstable equilibrium fails stability. These distinctions also apply to
continuous-time models.
\end{definition}
The sources often use ``stable'' to mean asymptotically stable. Here a
convergence conclusion is stated explicitly.

Write $x_t=x^*+y_t$. Since $f(x^*)=x^*$,
\[
y_{t+1}=f(x^*+y_t)-f(x^*)=f'(x^*)y_t+o(y_t).
\]
Discarding the smaller terms gives the \emph{linearised} recurrence
\[
y_{t+1}=\mu y_t,\qquad y_t=\mu^t y_0,\qquad \mu=f'(x^*).
\]
This explicit expression solves the linear approximation, not in general
the full nonlinear recurrence.

\begin{proposition}[Derivative test for a map]
For a continuously differentiable map near $x^*$:
\begin{itemize}
\item $|f'(x^*)|<1$ implies local asymptotic stability;
\item $|f'(x^*)|>1$ implies instability;
\item $|f'(x^*)|=1$ leaves the linear test inconclusive.
\end{itemize}
\end{proposition}
\begin{proof}[Explanation]
If $|f'(x^*)|<1$, continuity gives a neighbourhood with $|f'|\leq q<1$.
The mean value theorem yields $|f(x)-x^*|\leq q|x-x^*|$, so sufficiently
nearby iterates remain in the neighbourhood and converge geometrically.
If $|f'(x^*)|>1$, the same argument gives a local expansion factor greater
than one, so a nonzero perturbation eventually leaves a fixed small
neighbourhood. At equality, higher-order terms may decide the behaviour.
\end{proof}
For $0<\mu<1$, the linearised approach keeps the perturbation's sign.
For $-1<\mu<0$, it alternates sides while shrinking. At $\mu=0$,
higher-order terms determine the detailed approach, but the equilibrium
is still locally asymptotically stable.

\begin{example}[Logistic fixed-point stability]
Here $f'(x)=r(1-2x)$. At zero, $\mu=r$, so zero attracts for $0<r<1$
and is unstable for $r>1$. At $x^*=1-1/r$,
\[
\mu=r\left(1-2+\frac2r\right)=2-r.
\]
Thus the positive equilibrium attracts when $1<r<3$, and is unstable for
$r>3$. The approach is direct for $1<r<2$ and alternating for $2<r<3$.
The linear test alone does not settle $r=1$ or $r=3$.
\end{example}

\section{Parameters and bifurcations}
\begin{definition}[Bifurcation]
A bifurcation is a qualitative change in the dynamics as a parameter varies,
for example creation of equilibria or a change in their stability.
A bifurcation diagram plots equilibria or periodic states against a parameter;
solid and dashed lines often indicate attracting and unstable branches.
\end{definition}

\begin{example}[A parameter-dependent quadratic map]
For $x_{t+1}=x_t^2+x_t-a$, fixed points satisfy $(x^*)^2=a$.
For $a<0$ there are none; for $a=0$ there is one; for $a>0$ there are two:
$x^*=\pm\sqrt a$. Since $f'(x)=1+2x$,
\[
f'(\sqrt a)=1+2\sqrt a>1,\qquad
f'(-\sqrt a)=1-2\sqrt a.
\]
The positive branch is unstable for every $a>0$. The negative branch
attracts precisely in the range given by the strict derivative test,
$|1-2\sqrt a|<1$, that is $0<a<1$. It is unstable for $a>1$.
At $a=0$ and $a=1$, the multiplier equals $1$ or $-1$ respectively,
so that test is inconclusive at the parameter value itself.
These are the bifurcation values discussed in the sources.
\end{example}

\section{Period-two cycles and the route to chaos}
\begin{definition}[Period-two cycle]
Two distinct values $x_-,x_+$ form a period-two cycle if
$f(x_-)=x_+$ and $f(x_+)=x_-$. They satisfy
$f^{\circ2}(x)=x$, where $f^{\circ2}=f\circ f$ denotes composition, not
squaring the value $f(x)$.
\end{definition}
A fixed point of $f$ is also fixed by $f^{\circ2}$, so those old solutions
must be removed when looking for a genuine period-two cycle.

For the logistic map,
\[
f^{\circ2}(x)=r^2x(1-x)\bigl(1-rx(1-x)\bigr).
\]
Factoring $f^{\circ2}(x)-x=0$ gives the known roots $0,1-1/r$ and the
additional quadratic
\[
x^2-\left(1+\frac1r\right)x+\frac1r+\frac1{r^2}=0.
\]
Its two roots are
\begin{equation}\label{eq:applied-two-cycle}
x_\pm=\frac{r+1\pm\sqrt{r^2-2r-3}}{2r}.
\end{equation}
For $r>0$, they form a genuine distinct positive two-cycle when $r>3$.
At $r=3$ they merge with the fixed point; for the population interval,
we restrict attention to $r\leq4$.

\subsection{Stability of a cycle}
Let $g=f^{\circ2}$. By the chain rule, both cycle points have multiplier
\[
g'(x_-)=f'(x_+)f'(x_-)=g'(x_+).
\]
Using the sum and product of the two roots,
\begin{align*}
g'(x_\pm)
&=r^2(1-2x_-)(1-2x_+)\\
&=r^2\left[1-2\left(1+\frac1r\right)
+4\left(\frac1r+\frac1{r^2}\right)\right]\\
&=5-(r-1)^2.
\end{align*}
Hence the cycle attracts for
\[
|5-(r-1)^2|<1\quad\Longleftrightarrow\quad3<r<1+\sqrt6.
\]
A cycle can be attracting even though neither point is a fixed point of $f$:
we test the product of derivatives over a full cycle.

\subsection{Period doubling and chaotic behaviour}
At $r=3$, the fixed point loses attraction and the attracting two-cycle
appears. At $r=1+\sqrt6\approx3.44949$, the cycle multiplier reaches $-1$
and a four-cycle subsequently attracts. Further doublings give periods
$8,16,\ldots$, accumulating near $r\approx3.56995$.

Some larger parameters give bounded, aperiodic trajectories sensitive to
initial conditions: the source's introductory description of chaos.
Periodic windows also occur, including period three near
$1+\sqrt8\approx3.82843$; not every parameter beyond accumulation is chaotic.
